
In this phase, we elevate the representation from pixel-level trajectories to a semantic Knowledge Graph $\mathcal{G} = (\mathcal{V}, \mathcal{E})$. Each spatiotemporal tubelet $\mathcal{T}_k$ derived in the previous step is instantiated as a unique entity node $e_k \in \mathcal{V}$. The structure of this component is shown in the Fig. \ref{fig:level2}.
% 在此阶段，我们将表示从像素级轨迹提升到语义知识图谱 $\mathcal{G} = (\mathcal{V}, \mathcal{E})$。在前一步中导出的每个时空管 $\mathcal{T}_k$ 都被实例化为唯一的实体节点 $e_k \in \mathcal{V}$。

\textbf{Entity Relationship Modeling.}
% 实体关系建模。
To discover latent relationships between entities (e.g., \textit{Instrument A cuts Tissue B}), we evaluate pairwise affinities across the temporal dimension. We employ Video-CLIP \cite{videoclip}to encode the aggregated video segments of each tubelet, yielding global semantic embeddings. The relationship extraction logic considers both \textbf{long-term semantic interaction} and \textbf{instantaneous spatial proximity}.
% 为了发现实体间的潜在关系（例如，仪器 A 切割组织 B），我们评估了时间维度上的成对亲和力。我们采用 Video-CLIP 对每个管状结构的聚合视频片段进行编码，从而得到全局语义嵌入。关系抽取逻辑同时考虑了长期语义交互和瞬时空间邻近性。

We define the \textbf{Spatio-Temporal Overlap Score} ($S_{st}$) to capture physical interactions (e.g., touching, cutting) by averaging the GIoU over the co-existent frames $T_{overlap}$:
% 我们定义时空重叠分数（$S_{st}$）来捕捉物理交互（例如，触摸、切割），方法是对共存帧 $T_{overlap}$ 上的 GIoU 进行平均：
\begin{equation}
    S_{st}(\mathcal{T}_k, \mathcal{T}_q) = \frac{1}{|T_{overlap}|} \sum_{t \in T_{overlap}} \text{GIoU}(b_k^{(t)}, b_q^{(t)})
\end{equation}

Simultaneously, to capture correlated motion patterns (e.g., an instrument following a specific path relative to an organ), we calculate the \textbf{Trajectory Similarity} ($S_{traj}$) using Dynamic Time Warping (DTW). We convert the DTW distance into a similarity score via an exponential kernel:
% 同时，为了捕捉相关的运动模式（例如，器具相对于器官沿着特定路径运动），我们使用动态时间规整（DTW）计算轨迹相似度（$S_{traj}$）。我们通过指数核将DTW距离转换为相似度得分：
\begin{equation}
    S_{traj}(\mathcal{T}_k, \mathcal{T}_q) = \exp\left(-\gamma \cdot \text{DTW}(\mathbf{c}_k, \mathbf{c}_q)\right)
\end{equation}
where $\mathbf{c}$ denotes the sequence of centroid coordinates of the bounding boxes, and $\gamma$ is a decay hyperparameter controlling sensitivity.
% 其中 $\mathbf{c}$ 表示边界框的质心坐标序列，$\gamma$ 是控制灵敏度的衰减超参数。

The \textbf{Connection Priority Score} is computed by fusing the semantic cosine similarity ($S_{sem}$) with the motion metrics:
% 连接优先级得分是通过将语义余弦相似度（$S_{sem}$）与运动度量融合计算得出的：
\begin{equation}
    P_{edge}(\mathcal{T}_k, \mathcal{T}_q) = \sigma(S_{sem}) \cdot \left( \beta S_{st} + (1-\beta) S_{traj} \right)
\end{equation}
Pairs with $P_{edge}$ exceeding a confidence threshold are candidate edges. To resolve the specific semantic predicate of the edge (e.g., \textit{grasps}, \textit{cauterizes}), we feed the visual crops and metadata of the interacting tubelets into an MLLM (Qwen3-VL \cite{qwen3vl}). The MLLM acts as a neuro-symbolic reasoned, assigning explicit relationship labels to edges, thereby completing the construction of $\mathcal{G}$.
% 置信度超过阈值的 $P_{edge}$ 对被视为候选边。为了确定边的具体语义谓词（例如，\textit{抓取}、\textit{烧灼}），我们将相互作用的管状结构的视觉裁剪图和元数据输入到 MLLM（Qwen3-VL）中。MLLM 充当神经符号推理器，为边分配明确的关系标签，从而完成 $\mathcal{G}$ 的构建。

